\clearpage
\VerdierPhysicalPageAnchor{0262}%
\begingroup
\small
\setlength{\parindent}{0pt}%
\setlength{\parskip}{0pt}%
\setlength{\emergencystretch}{2em}%
\newcommand{\VerdierTermRow}[3]{%
  \VerdierSemanticAnchor{#1}%
  \par\noindent\hangindent=1.5em\hangafter=1
  #2。\enspace\textup{#3。}\par\smallskip}

\VerdierTermRow{idxt-p262-r01}{$\mathcal{F}\ell_n(\widetilde{\mathbb Z})$ の有限対象}{第 II 章、4.3.1}
\VerdierTermRow{idxt-p262-r02}{$\operatorname{comp}^{n}(C)$ の総和可能な対象}{第 I 章、2.1.7}
\VerdierTermRow{idxt-p262-r03}{$\displaystyle\prod_{i\in[n]}\operatorname{comp}(C_i)$ の総和可能な対象}{第 I 章、2.4.2}
\VerdierTermRow{idxt-p262-r04}{スペクトル対象}{第 II 章、4.3.2}
\VerdierTermRow{idxt-p262-r05}{$D(\mathcal A)$ に値をとるスペクトル対象、第 $m$ 次でフィルター付けられた複体により定まるもの}{第 III 章、4.1.8}
\VerdierTermRow{idxt-p262-r06}{$K(\mathcal A)$ に値をとるスペクトル対象、第 $m$ 次でフィルター付けられた複体により定まるもの}{第 III 章、4.1.7}
\VerdierTermRow{idxt-p262-r07}{加法圏の複体の圏に値をとる関手に付随するスペクトル対象}{第 II 章、4.1.3}
\VerdierTermRow{idxt-p262-r08}{Cartan--Eilenberg スペクトル対象}{第 III 章、4.6.7}
\VerdierTermRow{idxt-p262-r09}{アーベル圏に値をとる $\mathcal J$ 型スペクトル対象}{第 II 章、4.1.4}
\VerdierTermRow{idxt-p262-r10}{三角圏に値をとる $\mathcal J$ 型スペクトル対象}{第 II 章、4.1.2}
\VerdierTermRow{idxt-p262-r11}{定常スペクトル対象}{第 II 章、4.4.2}
\VerdierTermRow{idxt-p262-r12}{加法圏に値をとる関手に対する符号群の標準作用}{第 I 章、1.4.3}
\VerdierTermRow{idxt-p262-r13}{$\Pi$-拡張（関手 $\int_\varphi$）}{第 I 章、2.2.3}
\VerdierTermRow{idxt-p262-r14}{複体のホモトピー圏への関手の $\Pi$-拡張}{第 I 章、2.5.10}
\VerdierTermRow{idxt-p262-r15}{多加法関手の $\Pi$-拡張}{第 I 章、2.4.2}
\VerdierTermRow{idxt-p262-r16}{複体のホモトピー圏への関手の射の $\Pi$-拡張}{第 I 章、2.5.10}
\VerdierTermRow{idxt-p262-r17}{複体に付随する $D(\mathcal A)$ に値をとる第一スペクトル対象}{第 III 章、4.1.5}
\VerdierTermRow{idxt-p262-r18}{複体に付随する $K(\mathcal A)$ に値をとる第一スペクトル対象}{第 III 章、4.1.4}
\VerdierTermRow{idxt-p262-r19}{複体に付随する第一スペクトル対象}{第 III 章、4.1.5}
\VerdierTermRow{idxt-p262-r20}{$D(\mathcal A)$ に値をとる第一標準スペクトル対象}{第 III 章、4.1.5}
\VerdierTermRow{idxt-p262-r21}{$K(\mathcal A)$ に値をとる第一標準スペクトル対象}{第 III 章、4.1.4}
\endgroup
